\chapter{Security}
\label{ch:security}

\section{Method and scope}

The review covers the repository's own source, templates, build scripts and CI
configuration. Findings are stated in a fixed format: finding, severity, evidence,
impact and recommendation. Severity reflects impact in the product's real deployment
model --- a locally executed developer tool that writes files into the user's
project and may fetch content from a configurable registry.

\section{Secrets}

\begin{longtable}{@{}L{1.3cm}L{3.4cm}L{1.9cm}L{6.1cm}L{1.5cm}@{}}
\caption{Secret-related findings}\\
\toprule
\textbf{ID} & \textbf{Finding} & \textbf{Severity} & \textbf{Evidence} & \textbf{Status} \\
\midrule
\endfirsthead
\toprule
\textbf{ID} & \textbf{Finding} & \textbf{Severity} & \textbf{Evidence} & \textbf{Status} \\
\midrule
\endhead
SEC-01 & No secrets are committed. A recursive search for \texttt{.env} files found
none, and no credential literals appear in source, templates or workflows. & --- &
Absence verified across the tree. & PASS \\
SEC-02 & CI uses repository secrets for publishing
(\texttt{NPM\_TOKEN}) and the workflow default token
(\texttt{GITHUB\_TOKEN}), both referenced indirectly rather than inlined. &
--- & \texttt{publish.yml}, \texttt{deploy-registry.yml}. & PASS \\
SEC-03 & The CLI has no mechanism to store credentials; it therefore cannot leak
them. & --- & \texttt{registry.ts} performs unauthenticated GET only. & PASS \\
\bottomrule
\end{longtable}

\section{Injection}

\begin{longtable}{@{}L{1.3cm}L{3.6cm}L{1.9cm}L{5.8cm}L{1.6cm}@{}}
\caption{Injection findings}\\
\toprule
\textbf{ID} & \textbf{Finding} & \textbf{Severity} & \textbf{Evidence} & \textbf{Status} \\
\midrule
\endfirsthead
\toprule
\textbf{ID} & \textbf{Finding} & \textbf{Severity} & \textbf{Evidence} & \textbf{Status} \\
\midrule
\endhead
SEC-04 & \textbf{Shell command construction from unvalidated user input} in
\texttt{init}. The project name is interpolated directly into a shell string that
is executed with inherited stdio. A name containing shell metacharacters or quote
breakout would be interpreted by the shell. & High &
\texttt{packages/cli/src/commands/init.ts}, lines 20--23. & OPEN \\
SEC-05 & \textbf{Dependency names are interpolated into an executed shell
command.} Recipe \texttt{dependencies} strings are joined and passed to
\texttt{execSync}. With a bundled catalog this is trusted data; with a remote
registry, recipe content becomes an execution vector (see SEC-08). & High &
\texttt{packages/cli/src/utils/installDeps.ts}, lines 25 and 35. & OPEN \\
SEC-06 & No SQL, NoSQL or ORM access exists, so SQL injection is
\textbf{NOT APPLICABLE}. & --- & No database in the project. & N/A \\
SEC-07 & Template rendering is not a code-execution path: Handlebars is used in
default (non-helper-registering) mode with only the \texttt{eq} and \texttt{neq}
helpers added, and generated output is written to files rather than evaluated. & Low &
\texttt{packages/core/src/generate.ts}, lines 5--6. & ACCEPTED \\
\bottomrule
\end{longtable}

\subsection*{SEC-04 / SEC-05 --- details and recommendation}

\begin{lstlisting}[language=JavaScript]
// packages/cli/src/commands/init.ts (excerpt)
const nameArg = name ? `${name}` : ".";
const cmd = `npx create-expo-app@latest ${nameArg} --template ${template} --yes`;
execSync(cmd, { cwd: process.cwd(), stdio: "inherit" });
\end{lstlisting}

Both \texttt{name} and \texttt{template} come from argv. Neither is validated or
quoted.

\textbf{Impact:} a crafted project name (for example one containing
\texttt{;}, \texttt{\&\&} or a backtick) executes arbitrary commands with the user's
privileges. The risk is partly self-inflicted (the user types their own name), but
the command may also be supplied by a script, CI job or documentation snippet.

\textbf{Recommendation:} use \texttt{execFileSync} (or
\texttt{spawnSync}) with an argument array so no shell is involved, and validate
the project name against a conservative pattern
(\texttt{\^{}[A-Za-z0-9.\_-]+\$}). Apply the same treatment to the dependency
installer, and validate dependency strings against npm package-name grammar.

\section{Trust boundary: the registry}

\begin{longtable}{@{}L{1.3cm}L{3.6cm}L{1.9cm}L{5.8cm}L{1.6cm}@{}}
\caption{Remote-content findings}\\
\toprule
\textbf{ID} & \textbf{Finding} & \textbf{Severity} & \textbf{Evidence} & \textbf{Status} \\
\midrule
\endfirsthead
\toprule
\textbf{ID} & \textbf{Finding} & \textbf{Severity} & \textbf{Evidence} & \textbf{Status} \\
\midrule
\endhead
SEC-08 & \textbf{Remote recipes are rendered and written verbatim, with no
verification.} A configured registry can return any \texttt{files[].path} and any
template content; the client writes those paths into the user's project and
installs declared dependencies. & High &
\texttt{getRecipeById} (both clients) plus \texttt{writeFiles},
\texttt{installDeps}. & OPEN \\
SEC-09 & Path traversal is not validated on generated paths. A recipe path
containing \texttt{..} would resolve outside the project root because paths are
joined without containment checks. & High &
\texttt{writeFiles.ts} line 14; \texttt{generator.ts} line 215. & OPEN \\
SEC-10 & HTTPS is not enforced. A registry base of \texttt{http://} is accepted
without warning. & Medium & \texttt{getRegistryBase} accepts any string. & OPEN \\
SEC-11 & There is no integrity check (hash or signature) on fetched recipe
documents. & Medium & \texttt{fetchJSON} returns parsed JSON directly. & OPEN \\
\bottomrule
\end{longtable}

\textbf{Overall assessment.} Because the default configuration uses the embedded
catalog, the system is safe out of the box. The moment a remote registry is
configured --- which the product actively documents and supports --- the registry
becomes a trusted code-distribution channel with none of the controls a package
manager applies (no signature, no hash, no review, no path containment). This is
the single most important security consideration in the product.

\textbf{Recommendation (RECOMMENDED):}
\begin{itemize}
  \item reject or normalise any generated path that escapes the project root;
  \item validate dependency strings against the npm name grammar;
  \item warn or refuse when a non-HTTPS registry is configured;
  \item consider publishing a manifest of recipes with content hashes and
        verifying it client-side;
  \item document the trust implications of \texttt{--registry} prominently.
\end{itemize}

\section{Generated-code findings}

\begin{longtable}{@{}L{1.3cm}L{3.6cm}L{1.9cm}L{5.8cm}L{1.6cm}@{}}
\caption{Findings in generated screens}\\
\toprule
\textbf{ID} & \textbf{Finding} & \textbf{Severity} & \textbf{Evidence} & \textbf{Status} \\
\midrule
\endfirsthead
\toprule
\textbf{ID} & \textbf{Finding} & \textbf{Severity} & \textbf{Evidence} & \textbf{Status} \\
\midrule
\endhead
SEC-12 & \textbf{Credentials are logged to the console.} The generated login and
register hooks execute \texttt{console.log('Login attempt', \{ email, password
\})}, and the shadcn form logs the whole form payload. & Medium &
\texttt{hooks.use-login.hbs}; \texttt{components.login-form.shadcn.hbs}. & OPEN \\
SEC-13 & \textbf{Authentication is a placeholder.} Submission always succeeds after
a fixed delay and shows a success alert. If a developer ships without replacing it,
the app reports successful login for any credentials. & High &
\texttt{hooks.use-login.hbs}, \texttt{hooks.use-register.hbs}. & OPEN (documented) \\
SEC-14 & No secure storage is generated for any session token; the templates
perform no session persistence. This is a gap rather than a vulnerability, but it
means a developer following the scaffolding must add secure storage themselves. &
Low & Absence across auth templates. & INFO \\
SEC-15 & Sample data is embedded (statistics, product lists, notification lists,
placeholder names). No real data or personal data is present. & --- &
\texttt{packages/registry/templates/**}. & PASS \\
\bottomrule
\end{longtable}

\textbf{Recommendation:} remove the credential logging from the auth hook templates
(or gate it behind a development-only check), and state in the generated file's
leading comment and in user-facing documentation that the hook must be replaced
before the screen is functional.

\section{Input validation}

\begin{longtable}{@{}L{4.4cm}L{9.8cm}@{}}
\caption{Input validation posture}\\
\toprule
\textbf{Input} & \textbf{Validation} \\
\midrule
\endfirsthead
\toprule
\textbf{Input} & \textbf{Validation} \\
\midrule
\endhead
CLI option values (styling, router, ui, icons) & Checked against the recipe's
\texttt{supported} arrays; rejected with an error. \\
CLI screen id & Looked up in the catalog; rejected when absent. \\
CLI variant id & \emph{Not} validated. An unknown variant is silently ignored and the
default template is used. \\
CLI project name and template (\texttt{init}) & None. See SEC-04. \\
Registry URL & Trimmed and de-slash-normalised; otherwise unvalidated. \\
Remote response shape & Partially validated: index must be a non-empty array;
recipe must be an object with a truthy \texttt{id}. No field-level validation. \\
Recipe field values & Validated only by the build-time validator, not by clients. \\
Generated file paths & Not validated for containment. See SEC-09. \\
Recipe dependency strings & Not validated. See SEC-05. \\
\bottomrule
\end{longtable}

\section{Cross-site, CORS and web concerns}

\textbf{NOT APPLICABLE} for the product. There is no browser application, no HTML
rendering of remote content, no cookies and no session-bearing endpoints.

Related observations:
\begin{itemize}
  \item The development server sets
        \texttt{Access-Control-Allow-Origin: *}. This is appropriate for a
        localhost JSON fixture server but would be inappropriate if it were ever
        exposed publicly with sensitive content. It serves read-only static JSON, so
        the exposure is limited.
  \item The development server implements path-traversal protection by resolving
        the request against the registry root and rejecting escapes --- a good
        control, and notably absent from the file-writing path.
  \item The extension renders recipe descriptions in a
        \texttt{MarkdownString} in the completion list. VS Code sanitises
        documentation markdown by default; a malicious registry could still attempt
        content injection through that channel. Severity low, and only reachable
        with a configured untrusted registry.
\end{itemize}

\section{Dependency posture}

\begin{itemize}
  \item No dependency vulnerability scanning, audit step, Dependabot configuration
        or lockfile-integrity check exists in CI.
  \item Versions are declared with caret ranges (for example \texttt{commander
        \^{}12}, \texttt{handlebars \^{}4.7.8}, \texttt{esbuild \^{}0.21.5}), and a
        \texttt{pnpm-lock.yaml} is committed.
  \item \texttt{pnpm-workspace.yaml} explicitly allows build scripts for
        \texttt{esbuild} and \texttt{@vscode/vsce-sign}, and disables them for
        \texttt{keytar} --- a positive, deliberate posture-control measure.
\end{itemize}

\textbf{Status:} PARTIAL. \textbf{Recommendation:} add \texttt{pnpm audit} (or an
equivalent) to CI and enable dependency update automation.

\section{Positive security controls}

\begin{itemize}
  \item Network fetches are time-bounded (10\,s), so a hostile or hung endpoint
        cannot stall the client indefinitely.
  \item The development server defends against path traversal.
  \item Regeneration is diff-previewed and confirmation-gated by default.
  \item Primitives skip existing files, reducing accidental overwrites.
  \item Recipe validation runs before publication and checks structure and template
        existence.
  \item No telemetry means no unintended data egress.
  \item The offline-by-default design means most users never touch the trust
        boundary described in Section~\ref{ch:security}.
\end{itemize}

\textbf{Status:} IMPLEMENTED.

\section{Data integrity finding}

\begin{longtable}{@{}L{1.3cm}L{3.6cm}L{1.9cm}L{5.8cm}L{1.6cm}@{}}
\toprule
\textbf{ID} & \textbf{Finding} & \textbf{Severity} & \textbf{Evidence} & \textbf{Status} \\
\midrule
\endfirsthead
\toprule
\textbf{ID} & \textbf{Finding} & \textbf{Severity} & \textbf{Evidence} & \textbf{Status} \\
\midrule
\endhead
SEC-16 & \texttt{add} overwrites any existing file at a recipe path without warning
or backup (only primitives are skip-aware). & Medium &
\texttt{writeFiles.ts} line 16. & OPEN \\
SEC-17 & If dependency installation fails, files have already been written and the
manifest is not updated, leaving a project whose files and manifest disagree. &
Low & \texttt{add.ts}, lines 238--272 order. & OPEN \\
\bottomrule
\end{longtable}

\textbf{Recommendation (RECOMMENDED):} prompt before overwriting a file that was
not written by a previous generation (that is, one not listed in the manifest), and
record the manifest entry before installing dependencies (or roll back file writes
on install failure).

\section{Security summary}

\begin{table}[H]
\centering
\caption{Findings by severity}
\begin{tabular}{@{}L{2.6cm}L{11.6cm}@{}}
\toprule
\textbf{Severity} & \textbf{Findings} \\
\midrule
High & SEC-04 (init command injection), SEC-05 (dependency-name interpolation),
SEC-08 (unverified remote recipes), SEC-09 (no path containment), SEC-13 (placeholder
authentication). \\
Medium & SEC-10 (HTTPS not enforced), SEC-11 (no integrity check), SEC-12
(credential logging in generated code), SEC-16 (silent overwrite). \\
Low & SEC-07 (template execution surface --- accepted), SEC-14 (no secure storage
generated), SEC-17 (partial-failure consistency). \\
Pass & SEC-01, SEC-02, SEC-03, SEC-15 and the absence of SQL/CSRF/XSS surfaces. \\
\bottomrule
\end{tabular}
\end{table}

\textbf{Status:} RECOMMENDED (all remediations). None of the high-severity findings
is exploitable without a deliberate action by the user (supplying a crafted project
name, or configuring an untrusted registry), except SEC-12 and SEC-13, which are
properties of the code that the tool \emph{generates} and will reach users by
default.
